\documentclass[11pt]{amsart}

\usepackage{amsmath}
\usepackage{amssymb}
\usepackage{amsthm}
%\usepackage{psfig}

\begin{document}
\baselineskip=16pt
\textheight=8.5in
%\parindent=0pt 
\def\sk {\hskip .5cm}
\def\skv {\vskip .08cm}
\def\cos {\mbox{cos}}
\def\sin {\mbox{sin}}
\def\tan {\mbox{tan}}
\def\intl{\int\limits}
\def\lm{\lim\limits}
\newcommand{\frc}{\displaystyle\frac}
\def\xbf{{\mathbf x}}
\def\fbf{{\mathbf f}}
\def\gbf{{\mathbf g}}

\def\dbA{{\mathbb A}}
\def\dbB{{\mathbb B}}
\def\dbC{{\mathbb C}}
\def\dbD{{\mathbb D}}
\def\dbE{{\mathbb E}}
\def\dbF{{\mathbb F}}
\def\dbG{{\mathbb G}}
\def\dbH{{\mathbb H}}
\def\dbI{{\mathbb I}}
\def\dbJ{{\mathbb J}}
\def\dbK{{\mathbb K}}
\def\dbL{{\mathbb L}}
\def\dbM{{\mathbb M}}
\def\dbN{{\mathbb N}}
\def\dbO{{\mathbb O}}
\def\dbP{{\mathbb P}}
\def\dbQ{{\mathbb Q}}
\def\dbR{{\mathbb R}}
\def\dbS{{\mathbb S}}
\def\dbT{{\mathbb T}}
\def\dbU{{\mathbb U}}
\def\dbV{{\mathbb V}}
\def\dbW{{\mathbb W}}
\def\dbX{{\mathbb X}}
\def\dbY{{\mathbb Y}}
\def\dbZ{{\mathbb Z}}

\def\la{{\langle}}
\def\ra{{\rangle}}
\def\summ{{\sum\limits}}
\def\eps{{\varepsilon}}
\def\lam{{\lambda}}
\def\tr{{\rm tr}}
\def\char{{\rm char}}
\def\Mat{{\rm Mat}}
\def\Func{{\rm Func}}
\def\Span{{\rm Span}}
\def\Ker{{\rm Ker}}
\def\Im{{\rm Im}}



\bf\centerline{Homework \#3. Due on Thursday, Sep 17th, 11:59pm on Canvas}\rm
\vskip .1cm

\bf{Reading for this homework: }\rm Sections 1.6 and 2.1.

\bf{Plan for next week (Sep 14,16): }\rm  continuue with 2.1 (kernels and images of linear maps, rank-mullity theorem), 2.2 (matrix of a linear transformation), 2.3 (composition of linear maps and matrix multiplication) and start 2.4 (invertibility and isomorphisms).
\skv

\bf\centerline{Problems: }\rm
\skv

\bf{Problem 1: }\rm Let $V$ be a vector space over $\dbC$ with $\dim(V)=n$. As discussed in HW\#1,
we can also consider $V$ as a vector space over $\dbR$ (by restricting the scalars
from $\dbC$ to $\dbR$). Prove that as an $\dbR$-vector space $V$ has dimension $2n$.
{\bf Hint:} If $\{v_1,\ldots, v_n\}$ is a basis of $V$ as a vector space over $\dbC$,
how to get a basis of $V$ as a vector space over $\dbR$? Use the basic example
$V=\dbC$ to make a guess.
\skv

The next several problems in this homework deal with the notion of an internal direct sum of subspaces and the notion of a (vector space) complement. Let $V$ be a vector space
and $U,W$ subspaces of $V$. We say that $V$ is an (internal) direct sum of $U$ and $W$ if 
\begin{itemize}
\item[(1)] $V=U+W$ (that is, every element of $V$
can be written as $u+w$ for some $u\in U$ and $w\in W$) and 
\item[(2)] $U\cap W=\{0\}$. 
\end{itemize}
If $V$ is a direct sum of $U$ and $W$, we write
$V=U\oplus W$. 

Given a subspace $U$ of $V$, a \emph{(vector space) complement} of $U$ in $V$ is any subspace $W$ of $V$ such that $V=U\oplus W$
(unlike the set-theoretic complement, such $W$ is almost never unqiue).
\skv

{\bf Problem 2: } Let $V$ be a vector space, $U$ and $W$ subspaces of $V$, $\beta_U$ a basis for $U$ and $\beta_W$ a basis for $W$.
\begin{itemize}
\item[(a)] Prove that $V=U+W$ $\iff$ $\beta_U\cup\beta_W$ spans $V$;
\item[(b)] Prove that $U\cap W=\{0\}$ if and only if $\beta_U\cap \beta_W=\emptyset$ and $\beta_U\cup \beta_W$ is linearly independent.
\end{itemize}
Thus, $V=U\oplus W$ if and only if $\beta_U\cap \beta_W=\emptyset$ and $\beta_U\cup \beta_W$ is a basis of $V$.

\skv
{\bf Problem 3: } Let $V$ be a finite-dimensional vector space and $U$ a subspace of $V$. Use Problem~2 and a suitable theorem from class
to prove that
\begin{itemize}
\item[(i)] $U$ has a complement;
\item[(ii)] for any complement $W$ of $U$ we have $\dim(U)+\dim(W)=\dim(V)$. 
\end{itemize}
\skv
\skv
\bf{Problem 4: }\rm Let $V$ be a finite-dimensional vector space and $U$ and $W$ subspaces of $V$.
\begin{itemize}
\item[(a)] Prove that $\dim(U+W)=\dim(U)+\dim(W)-\dim(U\cap W)$. See Problem~29(a) in \S~1.6 for
a hint (this is a major generalization of 3(ii)).
\item[(b)] Prove that $V=U\oplus W$ if and only if $V=U+W$ and $\dim(U)+\dim(W)=\dim(V)$.
\item[(c)] Prove that $V=U\oplus W$ if and only if $U\cap W=\{0\}$ and $\dim(U)+\dim(W)=\dim(V)$.
\end{itemize}
\skv
\bf{Problem 5: }\rm Let $V=\dbR^2$.
\begin{itemize}
\item[(a)] Let $u=(1,2)$ and $U=\Span(u)$. Find infinitely many different complements for $U$.
\item[(b)] Let $U$ and $W$ be subspaces of $V$ with $\dim(U)=\dim(W)=1$. Prove that $W$ is a complement
of $U$ $\iff$ $W\neq U$.
\end{itemize}
\skv
{\bf Problem 6: } Let $F$ be a field, $V=F^n$ and $W=\{(x_1,\ldots, x_n)\in V: \sum\limits_{i=1}^n x_i=0\}$. Let $\{e_1,\ldots, e_n\}$
be the standard basis of $V$ (that is, $e_i$ is the vector whose $i^{\rm th}$ coordinate is $1$ and all other coordinates are $0$).
\begin{itemize}
\item[(1)] Let $U=\Span(e_1)$. Prove (directly from definition) that $U$ is a complement of $W$
\item[(2)] Let $S=\{e_i-e_1: 2\leq i\leq n\}$. Prove directly from definition that $S$ is both linearly independent and spanning for $W$
(and hence a basis for $V$). 
\item[(3)] Now use (1), an earlier problem in this HW and a theorem from class to explain why to show that $S$ from (b) is a basis
for $W$, it suffices to check that it is either linearly independent or spanning for $W$ (and there is no need to check both).
\end{itemize}
\skv
\bf{Problem 7: }\rm For each of the following maps $T$ do the following:
Prove that $T$ is linear,  find a basis for $\Ker(T)$ and $\Im(T)$ and then verify the assertion of the rank-nullity theorem. As usual, $\mathcal P_n(F)$ denotes the space of polynomials of degree $\leq n$ over 
a field $F$.
\begin{itemize}
\item[(a)] $T: \mathcal  P_6(\dbR)\to \mathcal  P_6(\dbR)$ given by $T(f(x))=f'(x)$
\item[(b)] $T: \mathcal  P_6(\dbZ_3)\to \mathcal  P_6(\dbZ_3)$ given by $T(f(x))=f'(x)$
(where as before $\dbZ_3$ is the field of congruence classes mod $3$).
{\bf WARNING:} The answer in (b) is different from the answer in (a).
\item[(c)] $T: \mathcal  P_3(\dbR)\to \dbR$ given by $T(f(x))=f(2)$, that is,
$T$ is the evaluation map at $x=2$.
\item[(d) ] $T: \mathcal  P_3(\dbR)\to \mathcal  P_4(\dbR)$ given by $T(f(x))=(x+1)p(x)$,
that is, $T$ is the multiplication by $x+1$.
\end{itemize}
\skv
\bf{Problem 8: }\rm Let $f:V\to W$ be a linear map, let $S=\{s_1,\ldots, s_n\}$ a finite subset of $V$ and $f(S)=\{f(s_1),\ldots, f(s_n)\}$.
Determine if the following 
statements are true (always) or false (in at least one case). If the statement is true, prove it; if false, give a specific counterexample.
\begin{itemize}
\item[(a)] If $S$ spans $V$, then $f(S)$ spans $\Im(f)$.
\item[(b)] If $S$ is linearly independent and the vectors $f(s_1),\ldots, f(s_n)$ are distinct, then $f(S)$ is linearly independent.
\item[(c)] If $f(S)$ spans $\Im(f)$, then $S$ spans $V$.
\item[(d)] If  the vectors $f(s_1),\ldots, f(s_n)$ are distinct and $f(S)$ is linearly independent, then $S$ is linearly independent.
\end{itemize}
\end{document}
